\section{Tecnologias e Organização}
\label{sec:tecnologias}

\begin{table}[htbp]
  \centering
  \caption{Ferramentas usadas no banco de dados}
  \label{tab:ferramentas}
  \begin{tabularx}{\textwidth}{@{}L{0.6}L{1.4}@{}}
    \toprule
    \poliCabTabela{Ferramenta} & \poliCabTabela{Papel} \\
    \midrule
    PostgreSQL 16 & Sistema gerenciador do banco, local (Docker) e em produção \\
    Flyway & Cria e altera o esquema com migrations em SQL, numeradas pela data e hora \\
    \campo{pg} (node-postgres) & Driver usado pela API para executar o SQL dos repositórios \\
    Prisma & Client tipado, gerado a partir do banco, usado no seed e nos testes \\
    Better Auth & Biblioteca de autenticação; usa as tabelas \campo{user}, \campo{session}, \campo{account} e \campo{verification} \\
    D2 & Desenha o diagrama do banco a partir do próprio Postgres, a cada push \\
    \bottomrule
  \end{tabularx}
\end{table}

A divisão de responsabilidades é estrita: \textbf{só o Flyway altera o
esquema}. O Prisma apenas espelha o banco (\campo{prisma db pull}), e o Better
Auth nunca cria tabelas sozinho. Quando a biblioteca precisa de uma coluna
nova, a equipe escreve uma migration com as convenções do projeto. Assim, há
um único histórico de mudanças, e qualquer máquina chega ao mesmo esquema
aplicando as migrations em ordem.

A API segue a arquitetura de portas e adaptadores. As regras de negócio ficam
em \campo{api/domain}, que declara \emph{portas} (interfaces como
\campo{PlaylistRepository}) sem conhecer o banco. As implementações em SQL
ficam em \campo{api/adapters/db}, uma por porta, como
\campo{PostgresPlaylistRepository}. Isso permite testar os casos de uso com
repositórios em memória e testar os adaptadores contra um Postgres real.
